\section{Entanglement Entropy}

\begin{parag}{Recap}
    Let do a quick recap about Von Neumann entropy. Given a system described by a density matrix $\rho \in \mathcal{H}$ which is Hermitian $\rho = \rho^{\dagger}$ and positive semi definite, $\rho \preceq 0$ with $\text{tr}\left(\rho\right)= 1$. This generalized the concept of pure state $\ket{\psi}$ to \important{mixed} state.\\
	$ $\\
	For a one-qubit system, one can always write the density matrix as 
	\begin{align*} 
		\mathbb{C}^2 \ni \rho = \frac{1}{2}\left(I + \vec{a}\cdot \vec{\sigma}\right)
	\end{align*}
	With $\vec{a} \in \mathbb{R}^3$ and further more $\vec{a} \in$ Ball of radius 1.\\
	$ $\\
	One can then define the Von Neumann entropy as 
	\begin{align*} 
		S\left(\rho\right) = - \text{Tr}\left(\rho \ln \rho\right)  =  - \sum_{i= 1}^{D} \lambda_i \ln \lambda_i
	\end{align*}
	With $\lambda_i$ being the eigenvalues of $\rho$ (as $\rho$ is PSD, we have that those are nonnegative). For a one qubit system, we had found 
	\begin{align*} 
		S\left(\rho\right) = - \frac{\left(1 + \left\|a\right\|\right)}{2}\ln\left(\frac{1 + \left\|a\right\|}{2}\right) - \left(\frac{1 - \left\|a\right\|}{2}\right)\ln\left(\frac{1 - \left\|a\right\|}{2}\right) 
	\end{align*}
	With $\frac{1 - \left\|a\right\|^2}{2} = \det \rho$. One can see here that the entropy is maximized whenever $\left\|a\right\|= 0$ which gives $\ln 2$ and minimized whenever $\left\|a\right\|= 1$.
	
\end{parag}


Let us consider a bipartite system $\mathcal{H}_{\text{total}} = \mathcal{H}_A \otimes \mathcal{H}_B$. In that case for a pure state we can divide our system in two:
\begin{itemize}
    \item \textbf{product state}: one is able to factor the state such that $\ket{\psi} = \ket{\varphi_A} \otimes \ket{\varphi_B}$
    \item \textbf{entangled state}: $\ket{\psi}$ cannot be factorized into two kets.
\end{itemize}
For mixed state, one can define a sort of analogous  notion with :
\begin{itemize}
    \item \textbf{separable states}: 
		\begin{align*} 
			\rho =  \sum_{i= 1}^{k}\underbrace{w_i}_{\geq 0} \rho_A^{\left(i\right)} \otimes \rho_B^{\left(i\right)} 
		\end{align*}
		However this is not a nice generalization of the product state, for us to have a more correct generalization, one would also need convexity in the '\textit{mixture}' of our state. We would one to add the requirement $\sum_{i= 1}^{k}w_i = 1$.
		\item \textbf{entangled states}: non separable
\end{itemize}
 

\subsection{Quantify entanglement for pure states}
The goal is given a state

 \begin{align*} 
	\ket{\psi} \in \mathcal{H}_{\text{total}} =  \mathcal{H}_{A} \otimes \mathcal{H}_B
\end{align*}
With a density matrix assiociated to $\ket{\psi}$ defined as 
\begin{align*} 
	\rho_{AB} = \ket{\psi}\bra{\psi} 
\end{align*}
\begin{parag}{Definition of entanglement entropy}
	\begin{theoreme}
	Given a pure state $\ket{\psi} \in \mathcal{H}_{\text{total}} = \mathcal{H}_A \otimes \mathcal{H}_B$ with its associated matrix $\rho_{AB} = \ket{\psi}\bra{\psi}$ with its partial density matrices
	\begin{align*} 
		\rho_A &=  \text{Tr}_{\mathcal{H}_B}\ket{\psi}\bra{\psi}\\
		\rho_B &= \text{Tr}_{\mathcal{H}_A} \ket{\psi}\bra{\psi}
	\end{align*}
	we have the following equality
	\begin{align*} 
		S\left(\rho_A\right)  = S\left(\rho_B\right)
	\end{align*}
	\end{theoreme}



	\begin{definition}{(entanglement entropy)}
	Define the partial density matrix of $\rho_{AB}$ describing a pure state 
	\begin{align*} 
		\rho_A &=  \text{Tr}_{\mathcal{H}_B}\ket{\psi}\bra{\psi}\\
		\rho_B &= \text{Tr}_{\mathcal{H}_A} \ket{\psi}\bra{\psi}
	\end{align*}
	We define the entanglement entropy as 
	\begin{align*} 
		S\left(\rho_A\right) = S\left(\rho_B\right) 
	\end{align*}
	\end{definition}
\end{parag}

\begin{parag}{Example}
    Given our state being a product state 
	\begin{align*} 
		\ket{\psi} &= \ket{\varphi_A} \otimes \ket{\varphi_B} , \rho = \ket{\psi}\bra{\psi} = \left(\ket{\varphi_A} \otimes \ket{\varphi_B}\right)\left(\bra{\varphi_A} \otimes \bra{\varphi_B}\right)\\
		\rho_A &= \ket{\varphi_A}\bra{\varphi_A}\\
		\rho_B &= \ket{\varphi_B}\bra{\varphi_B}
	\end{align*}
	It is easy to see that 
	\begin{align*} 
		S\left(\rho_A\right) = S\left(\rho_B\right) = 0
	\end{align*}
	Let us try with the following state 
	\begin{align*} 
		\ket{B_{00}}  &= \frac{1}{\sqrt{2}}\left(\ket{0_A}\otimes\ket{0_B} + \ket{1_A}\otimes \ket{1_B}\right)\\
		\rho_{AB} = \ket{B_{00}}\bra{B_{00}}
	\end{align*}
	Here it is easy as always to see that $S\left(\rho_{AB}\right)= 0$ as our state is pure. Now the intresting question arises when we ask about the entangled entropy. To do so we need to compute 
	\begin{align*} 
		\rho_A &= \frac{1}{2}\text{Tr}_{\mathcal{H}_B}\left(\ket{0_A} \otimes \ket{0_B} + \ket{1_A} \otimes \ket{1_B}\right)\left(\bra{0_A} \otimes \bra{0_B} + \bra{1_A} \otimes \bra{1_B}\right)\\
		&= \frac{1}{2}\text{Tr}_{\mathcal{H}_B}[\ket{0_A}\bra{0_A} \otimes \ket{0_B}\bra{0_B} + \ket{0_A}\bra{1_A} \otimes \ket{0_B}\bra{1_B}\\
		&\mathspace \mathspace+ \ket{1_A}\bra{0_A} \otimes \ket{1_B}\bra{0_B} + \ket{1_A}\bra{1_A} \otimes \ket{1_B}\bra{1_B}]\\
		&= \frac{1}{2} \left(\ket{0_A}\bra{0_A} + \ket{1_A}\bra{1_A}\right)
	\end{align*}
	Here we found that 
	\begin{align*} 
		\rho_A = \frac{1}{2}I_A 
	\end{align*}
	The entropy here gives 
	\begin{align*} 
		S\left(\rho_A\right) = \log 2 = S\left(\rho_B\right)
	\end{align*}
\end{parag}
\begin{framedremark}
Remark that this result is very '\textit{non-classical}' in the sense that a state that is totally '\textit{ordered}' becomes maximally unordered when looking at a smaller part of the system.\\
$ $\\
Classically, when working with Shannon entropy, we would have for two random variable $\left(X, y\right)$ the following inequality 
\begin{align*} 
	H\left(X, Y\right) \geq H\left(X\right) \left(H\left(Y\right)\right)
\end{align*}
This is one of the main difference between classical information theory and quantum information theory.
\end{framedremark}
\begin{parag}{Araki-Lieb inequality}
	\begin{theoreme}{(Araki-Lieb inequality)}
    \begin{align*} 
    	\left|S\left(\rho_A\right) - S\left(\rho_B\right)\right| \leq S\left(\rho_{AB}\right)  \leq S\left(\rho_A\right) + S\left(\rho_B\right)
    \end{align*}
    \end{theoreme}
	Which we will prove next lecture.
\end{parag}


\paragraph{Schmidt Theorem}%
\label{par:Schmidt Theorem}

\begin{theoreme}{(Schmidt Theorem)}
Let $\ket{\psi} \in \mathcal{H}_A \otimes \mathcal{H}_B, \rho_{AB} =  \ket{\psi}\bra{\psi}$ with 
\begin{align*} 
	 \rho_A = \text{Tr}_{\mathcal{H}_B} \rho_{AB}, \mathspace\mathspace \rho_B =  \text{Tr}_{\mathcal{H}_A}\rho_{AB}
\end{align*}
\begin{enumerate}
    \item The non-zero eigenvalues of $\rho_A$ and $\rho_B$ are identical.
    \item As a consequence $S\left(\rho_A\right) = S\left(\rho_B\right)$
\end{enumerate}

\end{theoreme}
Before starting the proof let us do a recap on SVD in order to make our life easier
\begin{parag}{Math recap on SVD}
    Given a matrix $M$ with dimension $d_A \times d_B$, SVD tells us that we can always write our matrix as 
	\begin{align*} 
		M= U \Sigma V^{\dagger} 
	\end{align*}
	Where U and $V$ are unitary matrices and $\Sigma$ is a matrix of singular values.
\end{parag}


\begin{proof}
    Recall that our state $\ket{\psi} \in \mathcal{H}_A \otimes \mathcal{H}_B$ meaning that we can write our state using an orthonormal basis which is itself a tensor product of our two Hilbert spaces 
	\begin{align*}  
		=\left\{\ket{i}_A\right\}  \otimes \left\{\ket{j}_B\right\}
	\end{align*}
	Which means that we can rewrite our state as 
	\begin{align*} 
		\ket{\psi} = \sum_{i= 1, j= 1}^{d_A, d_B} \psi_{i, j} \ket{i}_A \otimes \ket{j}_B
	\end{align*}
	One can then define $\Psi = \left(\psi_{ij}\right)$  which we can also  decomposed using SVD giving 
	\begin{align*} 
		\Psi = U \Sigma V^{\dagger} 
	\end{align*}
	Let us also walk through our state 
	\begin{align*} 
		\rho_{AB} = \ket{\psi}\bra{\psi} = \sum_{i, j} \sum_{k, l} \psi_{i, j}\psi_{k, l}^{\ast}\left(\ket{i}_A \bra{k}_A\right) \otimes \left(\ket{j_B}\bra{l_B}\right)
	\end{align*}
	Now let us compute $\rho_A$, which here we require a sum over all element where we have $j= l$ :
	\begin{align*} 
		\rho_A &=  \sum_{i, j} \sum_{k} \psi_{i, j}\psi_{k, j}^{\ast} \ket{i_A}\bra{k_A} = \sum_{i, k}\left(\sum_{j}\psi_{i, j}\psi_{k, j}^{\ast}\right)\ket{i_A}\bra{k_A}
	\end{align*}
	Here we see that $\rho_a = \Psi \Psi^{\dagger}$. By the same reasoning we can find out that $\rho_B = \Psi^{\dagger}\Psi$ which using SVD decompositions implies that both share the same nonzero eigenvalues.
\end{proof}

